\documentclass[11pt]{article}
\usepackage{amsmath, amssymb, amsthm, bm, booktabs}
\usepackage{geometry}
\usepackage{hyperref}
\usepackage{url}

\geometry{margin=0.7in}

\usepackage{titlesec}
\titlespacing*{\section}{0pt}{8pt}{4pt}
\titlespacing*{\subsection}{0pt}{6pt}{3pt}

\setlength{\abovedisplayskip}{6pt}
\setlength{\belowdisplayskip}{6pt}
\setlength{\abovedisplayshortskip}{4pt}
\setlength{\belowdisplayshortskip}{4pt}

\usepackage{enumitem}
\setlist{noitemsep, topsep=4pt}

\linespread{0.95}

\title{Code Search over CodeSearchNet with BM25 and CodeBERT Reranking}
\author{Information Retrieval Course Project}
\date{\today}

\begin{document}

\maketitle

\begin{abstract}
This report describes a search engine for Python code. The system searches
functions from the Python part of CodeSearchNet using natural language queries.
We implemented two retrieval pipelines. The first pipeline uses only BM25 and a
SQLite inverted index. The second pipeline first retrieves BM25 candidates and
then reranks them with dense CodeBERT embeddings stored in a FAISS index. We also
implemented a Gradio demo for interactive comparison. In the final evaluation,
BM25 was stronger than the simple CodeBERT reranker. This result is important
because it shows that a neural model does not automatically improve retrieval
quality without fine-tuning or score calibration.
\end{abstract}

\section{Introduction}

Code search is a common task for developers. A user writes a natural language
request, for example ``download video by url'' or ``parse json file'', and the
system must return useful code. This task is different from normal document
search because code contains function names, identifiers, library calls, and
short comments. The same meaning can be written with different words, and this
makes search difficult.

The goal of this project was to build and evaluate a reproducible code search
pipeline for Python functions. We compare a classical sparse retrieval method,
BM25, with a two-stage method that uses BM25 plus dense CodeBERT reranking.

\section{Data Loading and Preprocessing}

The dataset is the Python subset of CodeSearchNet \cite{codesearchnet}.

For each record we keep the function name, docstring, code body, split name,
repository name, file path, and source URL. We build one document text from the
function name, docstring, and code. This document text is used by CodeBERT. For
BM25 we also build a lexical representation.

The BM25 tokenizer uses the following preprocessing:
\begin{itemize}
    \item extracts words, identifiers, and numbers;
    \item splits identifiers by underscores;
    \item splits camelCase identifiers;
    \item lowercases all tokens;
    \item removes records with missing function name or missing code.
\end{itemize}

The processed output contains document files, metadata files, and corpus
statistics.

The final processed corpus contains 457461 Python functions. Table
\ref{tab:data} shows the split sizes.

\begin{table}[h]
\centering
\begin{tabular}{lrrr}
\toprule
Split & Documents & Average chars & Average lexical tokens \\
\midrule
Train & 412178 & 1333.83 & 177.12 \\
Valid & 23107 & 1430.96 & 191.79 \\
Test & 22176 & 1355.84 & 180.81 \\
\bottomrule
\end{tabular}
\caption{Processed CodeSearchNet Python corpus.}
\label{tab:data}
\end{table}

\section{Indexing and Retrieval Using Only BM25}

BM25 is the lexical baseline. We implemented an inverted index in SQLite. The
index contains four main tables:
\begin{itemize}
    \item \texttt{documents}: document id, split, function name, and document length;
    \item \texttt{postings}: term, document id, and term frequency;
    \item \texttt{vocabulary}: document frequency, collection frequency, and inverted document frequency (IDF);
    \item \texttt{statistics}: global corpus statistics.
\end{itemize}
\\

The BM25 formula uses $k_1 = 1.5$ and $b = 0.75$. The IDF is computed as:
\[
IDF(t) = \log \left(1 + \frac{N - DF(t) + 0.5}{DF(t) + 0.5}\right),
\]
where $N$ is the number of documents and $DF(t)$ is the document frequency of
term $t$.

For a query, the system tokenizes the query with the same tokenizer, loads
postings lists from SQLite, computes BM25 scores, and returns the top documents.

The system uses test docstrings as queries. The paired test function is treated
as the only relevant document. 

\section{Indexing and Retrieval Using BM25 plus CodeBERT Dense Index}

The second pipeline is a two-stage system. It does not replace BM25. It first
uses BM25 to get a candidate set, and then it reranks this candidate set with
CodeBERT \cite{codebert}.

The system encodes each document text with \texttt{microsoft/codebert-base}. We
use mean pooling over token embeddings and then L2-normalize the result. The
embedding size is 768. The vectors are stored in a FAISS index with inner product
metric. Because the vectors are normalized, inner product is equal to cosine
similarity.

Then we evaluate this approach. For each query, BM25 returns 50 candidates.
Then CodeBERT computes cosine similarity between the query vector and candidate document vectors. The final
ranking is sorted by dense similarity. If two candidates have the same dense
score, BM25 score is used as a secondary score.

\section{Evaluation}

We evaluate both methods on the full test split with 22176 queries. The metrics
are MRR@10, Recall@10, Recall@50, and nDCG@10. Table \ref{tab:metrics} shows the
final results.

\begin{table}[h]
\centering
\begin{tabular}{lrrrr}
\toprule
Method & MRR@10 & Recall@10 & Recall@50 & nDCG@10 \\
\midrule
BM25 & 0.939881 & 0.983451 & 0.992109 & 0.950830 \\
BM25 + CodeBERT rerank & 0.195510 & 0.340458 & 0.992109 & 0.229310 \\
\bottomrule
\end{tabular}
\caption{Retrieval quality on the CodeSearchNet Python test split.}
\label{tab:metrics}
\end{table}

The result shows that BM25 is much better in this setup. This can look
unexpected, because CodeBERT is a neural model for code and natural language.
CodeBERT is based on the Transformer architecture and is pretrained on pairs of
natural language and source code. The pretrained model
\texttt{microsoft/codebert-base} was trained with CodeSearchNet-style data,
including documentation-code pairs. However, this does not mean that it will
automatically improve our reranking results. In our project, CodeBERT is used as
a zero-shot embedding model. We do not fine-tune it for our exact retrieval
objective. We only average token embeddings with mean pooling and compare vectors
with cosine similarity.

There are several reasons why this simple dense reranking gives worse results
than BM25.

First, the evaluation query is the original docstring of the target function.
This gives BM25 a strong advantage, because the docstring is also part of the
indexed document. There is often high lexical overlap between the query and the
correct document.

Second, the dense method can find semantically similar functions, but these
functions are not always the paired ground-truth function. For example, several
functions can implement the same general task, but the evaluation accepts only
one paired function as relevant. In this case, CodeBERT may move a similar but
wrong function above the exact target function, and the metric becomes lower.

Third, the dense stage only reranks the top 50 BM25 candidates. This explains
why Recall@50 is the same for both methods. The relevant document is in the same
candidate set, but the dense reranker often moves it lower inside the top 50.

This experiment gives a useful conclusion. A dense model can add semantic
matching, but it must be trained or combined carefully. For this current project,
BM25 is the best method by the measured ranking metrics.

\section{Gradio App for Demo}

The project includes a Gradio application.

The app has two tabs. The first tab is \textit{Search}. The user enters a
natural language query and sees two columns:
\begin{itemize}
    \item BM25 lexical results;
    \item BM25 + CodeBERT reranked results.
\end{itemize}

Each result shows the rank, function name, repository name, score, source link,
docstring, and code snippet. The user can change the number of displayed results
and the number of BM25 candidates for reranking. CodeBERT and FAISS are loaded
lazily on the first dense query, so the first dense query can be slower.

The second tab is \textit{Evaluation metrics}. It reads the precomputed JSON
files and shows the metric table, sample rankings, and raw JSON outputs. This
tab is useful for the screencast because it directly compares the two pipelines.

\section{Limitations and Future Work}

The main limitation is that CodeBERT is used without fine-tuning. Also, the
current reranker sorts by dense score only, so it can ignore very strong BM25
evidence. In future work, we would try three improvements:
\begin{itemize}
    \item fine-tune CodeBERT or another code retrieval model on CodeSearchNet;
    \item combine BM25 and dense scores with a weighted formula;
    \item evaluate on user-written queries, not only original docstrings.
\end{itemize}

\section{Conclusion}

We implemented a complete code search pipeline with data loading, preprocessing,
BM25 indexing, dense CodeBERT indexing, evaluation scripts, and a Gradio demo.
The main experimental result is that BM25 is the strongest method for this
evaluation setup. BM25 reaches MRR@10 of 0.939881, while BM25 + CodeBERT rerank
gets 0.195510. The result does not mean that dense retrieval is useless. It means
that simple pretrained embeddings and pure dense reranking are not enough for
this task. A better dense system needs fine-tuning or a hybrid score.

\begin{thebibliography}{9}

\bibitem{codesearchnet}
Husain, H., Wu, H., Gazit, T., Allamanis, M., and Brockschmidt, M.
CodeSearchNet Challenge: Evaluating the State of Semantic Code Search.
\textit{arXiv preprint arXiv:1909.09436}, 2019.
\url{https://arxiv.org/abs/1909.09436}

\bibitem{codebert}
Feng, Z., Guo, D., Tang, D., Duan, N., Feng, X., Gong, M., Shou, L., Qin, B.,
Liu, T., Jiang, D., and Zhou, M.
CodeBERT: A Pre-Trained Model for Programming and Natural Languages.
\textit{Findings of EMNLP}, 2020.
\url{https://arxiv.org/abs/2002.08155}

\bibitem{bm25}
Robertson, S. and Zaragoza, H.
The Probabilistic Relevance Framework: BM25 and Beyond.
\textit{Foundations and Trends in Information Retrieval}, 2009.

\end{thebibliography}

\end{document}
